%!TEX root = TQFTnotes.tex
%%%%%%%%%%%%%%%%%%%%%%%%%%%%%%
\chapter{Algebraic prerequisites: monoidal categories}\label{c:categories}

In this chapter we introduce some of the algebraic language
necessary for describing TQFTs, most importantly the notion
of a monoidal category. We only give a very brief introduction,
referring the reader to books \ocite{mac-lane}, \ocite{etingof-fusion}
for details.

%%%%%%%%%%%%%%%%%%
\section{Category theory basics}\label{s:cat-basics}
We assume that the reader is familiar with basic notions of categories and
functors. We will completely ignore all set-theoretic difficulties related
to the definition of a category; if necessary, the reader can assume that
all sets of objects, morphisms, etc in all the categories discussed are
subsets of a large enough universal set $U$.

For a category $\calC$, we will use the notation $X\in \calC$ to show that
$X$ is an object of $\calC$. For two objects $X,Y\in\calC$, we denote by
$\Hom_\calC(X,Y)$ the set of morphisms from $X$ to $Y$ (we will drop
the subscript $\calC$ when there is no ambiguity).

Given categories $\calC, \calD$, we denote by $\Fun(\calC, \calD)$ the
set of all functors $\calC\to \calD$.

Examples:
\begin{itemize}
  \item $\sets$: category of all sets.
  \item $\smoothman_n$: category of smooth manifolds of dimension $n$.
    There are different versions: oriented manifolds, framed manifolds, etc.
  \item $\plman_n$: category of PL manifolds of dimension $n$. Note that
    in dimensions $n\le 6$, every PL manifold has a canonical smooth structure
    and vice versa. In general, it is not so: every smooth manifold has
    a unique PL structure (Whitehead's theorem), but a PL manifold might
    not have a compatible smooth structure, or if it has one, it might not
    be unique (Milnor's exotic spheres). We will talk about PL manifolds in
    more detail later (see \chref{c:PL}).
  \item $\vect_\kk$: category of vector spaces over a field $\kk$.
  \item $R\mod$: category of left modules over an associative ring $R$.
\end{itemize}

One special case of a category, which we will use frequently, is a groupoid.
\begin{definition}\index{groupoid}
  A groupoid is a category in which all morphisms are invertible.
\end{definition}
\begin{example}
  Let $X$ be a topological space. Define its \emph{fundamental groupoid}\index{fundamental groupoid}
  $\pi(X)$ as the category whose objects are points in $X$ and morphisms
  are homotopy equivalence classes of paths between points, with the obvious
  composition.
\end{example}

One common theme in all our constructions is equality vs isomorphism. When
talking about points in a set (e.g., elements of a ring), we can require them to
be equal; for example, we can require $a+b=b+a$ for elements of a ring. However,
while it is technically possible to require two objects in a category to be
equal, it rarely makes sense to do so. For example, in the category of finite
sets, sets $A\times B$ and $B\times A$ are not equal. On the other hand,
requiring two objects to be isomorphic is typically too weak: e.g., any two sets
with the same number of elements are isomorphic.

The common way to resolve this is the following convention.
\begin{convention}\label{conv:canonical}
  A collection of objects $X_\al$ of category $\calC$, together with a
  collection of isomorphisms $f_{\be\al}\colon X_\al \isoto X_\be$, such that
  $f_{\al\be}f_{\be\ga}=f_{\al\ga}$, $f_{\al\al}=\id_{X_\al}$, is treated as a
  single object in $\calC$.
\end{convention}
This convention, though rarely stated explicitly, is actually widespread.
Virtually every math textbook treats sets $A\times B$ and $B\times A$ as the
same set, even though they are technically not the same object of category
$\sets$; from the point of view of the convention above, it is justified by
noting that we have a canonical isomorphism $A\times B\isoto B\times A\colon
(a,b)\mapsto (b,a)$.

Formally, this convention can be justified by constructing a category
$\widehat\calC$, whose objects are such collections $(X_\al, f_{\al\be})$, and
observing that categories $\calC$, $\widehat\calC$ are equivalent; see, e.g.,
\ocite{bakalov-kirillov}*{Lemma 1.1.12}.


%%%%%%%%%%%%%%
\section{Abelian categories}\label{s:abelian-cats}
\begin{definition}\label{d:abelian}\index{abelian category}
  An abelian category is a category in which for every pair of objects $X,Y$,
  the set of morphisms $\Hom(X,Y)$ is an abelian group, composition is
  additive, we have a zero object, the direct sum of two objects, and for every
  morphism, we have the notion of image and kernel satisfying the usual
  properties (see \ocite{etingof-fusion}*{Section~1.3} for details).
\end{definition}
In such a category, we can define the notion of a simple object (which has no
subobjects other than $0$ and itself), semisimple one (isomorphic to a direct
sum of simple objects), etc. We can also define the notions of short exact
sequences, Jordan--H\"older series (which in turn allows one to define the
notion of length of an object), injective and projective objects, and all
other usual constructions of homological algebra.

An abelian category is
\emph{$\kk$-linear}\index{category!kk-linear@$\kk$-linear} if all sets
$\Hom(X,Y)$ are vector spaces over a field $\kk$, and composition of morphisms
is linear over $\kk$.

\begin{example}
  Below are some examples of $\kk$-linear abelian categories.
  \begin{enumerate}
    \item $\vect_\kk$;
    \item category $\Rep_\kk(G)$ of representations of a group $G$;
    \item category $\Rep_\kk(\g)$ of representations of a Lie algebra $\g$;
    \item category $A\mod$ of modules over an associative algebra $A$
      over $\kk$.
  \end{enumerate}
\end{example}

Typically, one imposes some finiteness conditions for $\kk$-linear
categories. We will usually assume that all linear categories, unless
explicitly stated otherwise, are
\emph{locally finite}\index{category!locally finite}, i.e.,
\begin{itemize}
  \item All spaces $\Hom(X,Y)$ are finite-dimensional
  \item Every object has finite length
\end{itemize}
(see \ocite{etingof-fusion}*{Section 1.8}).

One can also impose stricter requirements, for example semisimplicity (every
short exact sequence splits), finiteness (there are finitely many isomorphism
classes of simple objects), and more. We will discuss it in more detail when it
becomes necessary.

One useful fact about $\kk$-linear categories is the following version of the
famous Schur's lemma in representation theory.

\begin{lemma}\label{l:schur}
  Let $\calA$ be a locally finite $\kk$-linear category.
  \begin{enumerate}
    \item If $X,Y\in \calA$ are simple objects, then any morphism
      $f\colon X\to Y$ is either zero or invertible.
    \item If $\kk$ is algebraically closed, then for any non-zero
      simple $X\in \calA$, $\Hom_\calA(X,X)=\kk$.
  \end{enumerate}
\end{lemma}
As with the usual Schur's lemma, the proof is immediate from the fact that the
only  finite-dimensional division algebra over an algebraically closed field
$\kk$ is $\kk$ itself (see, e.g., \ocite{pierce}*{Lemma 3.5}).

%%%%%%%%%%%%%%%%%%%%%%%%%%%%%%%
\section{Monoidal categories}\label{s:monoidal}
One of the most important notions is the notion of a monoidal category, which
informally is described as a category with an associative operation; the model
example is the category of vector spaces with the operation of tensor product.
In this section, we review this notion; as before, we refer the reader to
\ocite{etingof-fusion}*{Chapter 2}; \ocite{leinster}*{Section 1.2} for a more
in-depth discussion.

Before defining monoidal categories, let us look at simpler notion: a monoid.

A \emph{monoid}\index{monoid} is a set together with an associative
multiplication operation and a unit:
\begin{align*}
  &(a\cdot b)\cdot c = a\cdot (b\cdot c)\\
  &1\cdot a = a\cdot 1 = a
\end{align*}
(if we also add the requirement that every element has an inverse,
we get the definition of a group).

\begin{theorem}\label{t:coherence-monoids}
  Let $M$ be a monoid and $a_1, \dots, a_n\in M$. Consider all possible
  expressions obtained by writing $a_1, \dots, a_n$ (in this order), then
  inserting some copies of $1$ and then putting parentheses so that
  multiplications are just binary products, e.g., $(a_1 \cdot 1) \cdot (1 \cdot
  a_2) $. Then all of them are equal: if $X_1, X_2$ are two such expressions
  (with the same $a_1, \dots, a_n$), then $X_1=X_2$.
\end{theorem}
It may seem obvious, but in fact it is not quite trivial; essentially it is the
statement that if we construct a graph whose vertices correspond to
different expressions written above, and edges to application of associativity
and unit relations, then this graph is connected.

If we only include $a_1, \dots, a_n$ but no 1's, then such expressions are
naturally described by planar binary trees. For example, the bracketings $((a_1
a_2)a_3) a_4$ and $((a_1 a_2)(a_3 a_4))$ correspond to the trees below:

\begin{center}
  %%%%% ((12)3) 4
  \includegraphics{figures/fig01.svg}
  \qquad
  %%%%% ((12)(3 4)
  \includegraphics{figures/fig02.svg}
\end{center}
Adding ones requires a minor modification, adding appropriately
colored leaves to the tree.

For categories, the analogous structure is that of a monoidal category: a
category with a functor $\otimes \colon \calC\times \calC\to \calC$ and a unit
object $\one\in \calC$.

One can try to mimic the definition above.
\begin{definition}\label{d:strict_monoidal}
  A \emph{strict monoidal category}\index{monoidal category!strict} is a
  category $\calC$ together with a functor $\otimes \colon \calC\times \calC\to
  \calC$ and a unit object $\one\in \calC$ such that
  \begin{align*}
    &(A\otimes B)\otimes  C = A\otimes  (B\otimes C),\\
    &\one \otimes A  = A\otimes \one   = A.
  \end{align*}
\end{definition}

However, this definition is too strict: objects in a category are almost never
equal, just isomorphic. For example, if $V$ is a vector space over $\kk$, then
$\kk\otimes V\ne V$ (but is canonically isomorphic to $V$).

We can try to replace equality by isomorphism, but then it is too weak: e.g.,
any two vector spaces of the same dimension are isomorphic.

The proper way is to require the existence of a distinguished (canonical)
isomorphism which should be made part of the structure. Moreover, these
``canonical'' isomorphisms should themselves satisfy some compatibility
conditions, as suggested by Convention~\ref{conv:canonical}.

\begin{definition}\label{d:monoidal}\index{monoidal category}
  A monoidal category is a category $\calC$ together with the following data:
  \begin{itemize}
    \item An object $\one \in \calC$
    \item a functor $\otimes\colon \calC \times \calC \to \calC$
    \item functorial isomorphisms
      \begin{align*}
        &l_A\colon \one \otimes A \isoto A\\
        & r_A\colon A\otimes \one \isoto A\\
        &\al_{A,B,C}\colon (A\otimes B)\otimes C\isoto A\otimes (B\otimes C)
      \end{align*}
  \end{itemize}
  satisfying compatibility conditions below:
  \begin{enumerate}
    \item Triangle axiom: for any $A,B\in \calC$, the diagram below commutes
      \begin{center}\includegraphics{figures/fig03.svg}\end{center}
    \item Pentagon axiom: for any $A,B,C,D \in \calC$, the diagram below
      commutes
      \begin{center}\includegraphics{figures/fig04.svg}\end{center}
      (all arrows are given by the associativity isomorphism $\al$.)
%      \begin{equation}\label{e:pentagon}
%        \begin{tikzcd}[column sep=huge,row sep=large]
%          ((A\otimes B)\otimes C)\otimes D
%          \arrow[r,"\alpha_{A,B,C}\otimes \mathrm{id}_D"]
%          \arrow[d,"\alpha_{A\otimes B,C,D}"']
%          & (A\otimes (B\otimes C))\otimes D
%          \arrow[r,"\alpha_{A,B\otimes C,D}"]
%          & A\otimes((B\otimes C)\otimes D)
%          \arrow[d,"\mathrm{id}_A\otimes \alpha_{B,C,D}"] \\
%          (A\otimes B)\otimes(C\otimes D)
%          \arrow[rr,"\alpha_{A,B,C\otimes D}"']
%          && A\otimes(B\otimes(C\otimes D))
%        \end{tikzcd}
%      \end{equation}

  \end{enumerate}
\end{definition}

\begin{example}
The following are examples of monoidal categories.
  \begin{enumerate}
    \item Category $\sets$, with $\otimes$ given by disjoint union and
      $\one$ being the empty set
    \item Category $\sets$, with $\otimes$ given by the Cartesian product
      and $\one$ being a one-point set
    \item Categories $\smoothman$, $\plman$ with respect to disjoint union
    \item Category $\vect$ with the usual tensor product, $\one = \kk$
    \item Category of bimodules over a fixed ring $R$, with respect
      to $\otimes_R$, and $\one = R$

  \end{enumerate}
\end{example}


\begin{example}\label{x:vecg}
  Let $G$ be a group. Define the category $\vect_G$ of \emph{$G$-graded vector
  spaces}\index{G-graded vector spaces@$G$-graded vector spaces} as the category
  whose objects are finite-dimensional vector spaces $V$ together with
  decomposition
  $$
    V=\bigoplus_{g\in G} V_g.
  $$
  This is a monoidal category with tensor product given by
  $$
    (V\otimes W)_g = \bigoplus_{h\in G} V_h \otimes W_{h^{-1}g}
  $$
  and unit object given by $\kk$ considered as a $G$-graded vector space with
  grading equal to $1\in G$. The associativity isomorphism is trivial.
\end{example}



\begin{theorem}[Mac Lane's coherence theorem]\label{t:maclane}\par

  \noindent
  Let $\calC$ be a monoidal category. Let $F_1, F_2\colon \calC^n\to \calC$
  be two functors obtained by compositions of tensor product functor $\otimes$
  and functors $X\mapsto X\otimes \one$, $X\mapsto \one \otimes X$, e.g.,
  \begin{align*}
    F_1(A,B,C)&=(A\otimes B)\otimes (\one\otimes C),\\
    F_2(A,B,C)&=A\otimes (( B\otimes  C)\otimes \one).
  \end{align*}
  Then
  \begin{enumerate}
    \item There exists a functorial isomorphism $F_1\simeq F_2$, obtained
      by composition of associativity isomorphism $\al$, left and right
      unit isomorphisms $l,r$, and their inverses.
    \item Any two such isomorphisms $F_1\to F_2$ are equal.
  \end{enumerate}
\end{theorem}
In other words, if we construct a cell complex $X_n$ whose vertices are functors
$\calC^n\to \calC$ as in the statement of the theorem, edges correspond to
applications of the associativity and unit isomorphisms, and 2-cells correspond
to commutative diagrams given by the pentagon axiom, unit axioms, and
functoriality of $\al, l, r$, then this complex is simply-connected.

This is a natural generalization of coherence theorem for monoids
(\thref{t:coherence-monoids}); in this language, \thref{t:coherence-monoids}
is exactly claim (1) of the theorem, i.e., the statement that $X_n$
is connected.

\begin{remark}\label{r:stasheff}
  If we only consider functors obtained by composing the tensor product functor
  (thus, no units), the corresponding complex is a 2-skeleton of a convex
  polytope $K_n$, called the \emph{associahedron}\index{associahedron}, or the
  \emph{Stasheff
  polytope}\index{Stasheff polytope}. The vertices of $K_n$ are indexed by
  binary trees with $n$ leaves; thus, the number of vertices is the Catalan
  number $C_{n-1} = \frac{1}{n}\binom{2(n-1)}{n-1}$ (so $K_5$ has
  $C_4 = 14$ vertices).
  This polytope was introduced by Jim Stasheff in the 1960s. It can be realized
  as a convex polytope in Euclidean space; see, e.g.,
  \ocite{loday}.
\end{remark}

\begin{exercise}\label{ec:K5}
  Draw the associahedron $K_5$ (it has 14 vertices) and verify that it is
  simply-connected (more precisely, its 2-skeleton is topologically a sphere
  $S^2$). We will return to it later, when talking about monoidal 2-categories.
\end{exercise}

\begin{definition}\label{d:monoidal_functor}\index{monoidal functor}
  Let $\calC, \calD$ be monoidal categories. A monoidal functor
  $\calC\to \calD$ is a functor $F\colon\calC\to\calD$ together
  with functorial isomorphisms
  \begin{align*}
    J&\colon F(X\otimes Y)\to F(X)\otimes F(Y)\\
    \ph&\colon \one_\calD\to F(\one_\calC)
  \end{align*}
  satisfying suitable constraints (see \ocite{etingof-fusion}*{Section 2.4}).

  We say that categories $\calC, \calD$ are monoidally equivalent
  if there exist monoidal functors $F\colon \calC\to \calD$, $G\colon \calD\to
  \calC$ such that $F\circ G\simeq \id_\calC$, $G\circ F \simeq \id_\calD$.
\end{definition}


\begin{example}
  Consider the category of representations of a group $G$. Then the
  forgetful functor from this category to the category of vector
  spaces is a monoidal functor. More generally, given
  a subgroup $H\subset G$, the restriction functor
  $\Res^G_H\colon \Rep G\to\Rep H$ has an
  obvious monoidal structure.
\end{example}
\begin{example}
  If $G$ is a connected, simply-connected Lie group, then the category
  of finite-dimensional representations of $G$ is equivalent, as a
  monoidal category, to the category of finite-dimensional
  representations of its Lie
  algebra $\g$.
\end{example}

\begin{example}\label{x:vecg-twisted}
  Recall the category of $G$-graded vector spaces $\vect_G$ defined in
  \exref{x:vecg}. This category has a generalization. Namely, let
  $\om\colon G\times G\times G\to \kk^\times$ be a 3-cocycle, i.e., a
  function satisfying the condition
  $$
  \om(g_1g_2, g_3, g_4)\om (g_1, g_2, g_3g_4) = \om(g_1, g_2, g_3)
     \om(g_1, g_2g_3, g_4)\om (g_2, g_3, g_4).
  $$
  Consider the category $\vect^\om_G$ which coincides with $\vect_G$ as an
  abelian category; moreover, it has the same tensor product functor $\otimes$
  and unit object $\one=\kk_1$. However, the associativity isomorphism is
  different: namely, $\al \colon (V\otimes W)\otimes U \to V \otimes (W\otimes
  U)$ is given by
  $$
    \al((v \otimes w) \otimes u )= \om_{g,h,m} v \otimes (w\otimes u),
    \qquad v\in V_g, w\in W_h, u\in U_m.
  $$
  The cocycle condition immediately implies that so defined $\al$ satisfies the
  pentagon axiom.
\end{example}
\begin{exercise}
  Let $\om\colon G\times G\times G\to \kk^\times$ be a 3-cocycle. Assume that
  $\om$ is a coboundary, i.e., there exists a function $\be\colon G\times
  G\to \kk^\times$ such that $\om=d\be$, where the coboundary operator $d$
  is given by
  $$
    (d\be)(g_1, g_2, g_3)=
    \frac{\be(g_2, g_3)\,\be(g_1, g_2g_3)}{\be(g_1g_2, g_3)\,\be(g_1, g_2)}.
  $$
  Show that then category $\vect_G^\om$ is monoidally equivalent to $\vect_G$.

  More generally, categories $\vect_G^\om$ and $\vect_G^{\om'}$ are equivalent,
  with the equivalence preserving the $G$-graded dimension, if and only if
  $\om'=\om \cdot d\be$ for some $\be\colon G\times G\to \kk^\times$ (see
  \ocite{etingof-fusion}*{Section 2.6}).
\end{exercise}



We can restate Mac Lane's coherence theorem as follows.
\begin{theorem}\label{t:strictify}
  Any monoidal category is monoidally equivalent to a strict monoidal category.
\end{theorem}
If one knows the concept of a skeleton of a category, then Mac Lane's coherence
theorem can be understood as the construction of the skeleton by looking at the
isomorphism classes of objects in a monoidal category.






\begin{convention}\label{c:graphical-monoidal}
We will frequently use graphical presentation of morphisms
in a monoidal category, representing a morphism
$\ph \colon A_1\otimes \dots \otimes A_k\to B_1\otimes\dots \otimes B_m$
by a box with $k$ inputs at the top and $m$ outputs at the bottom:
\begin{figure}[h]
  \centering
  \includegraphics{figures/fig05.svg}
  \caption{Graphical notation for a morphism in a monoidal category.}
  \label{f:graphical-morphism}
\end{figure}
\end{convention}







%%%%%%%%%%%%%%%%%%%%%%%%%%%%%%%%%%%%%%%
\section{Symmetric monoidal categories}\label{s:symmetric-monoidal}
For the category of vector spaces, the tensor product is not only associative,
but also ``commutative'': one has natural isomorphisms $V\otimes W \simeq
W\otimes V$. The definition below gives a proper generalization of this to
arbitrary monoidal categories.

\begin{definition}\label{d:symmetric}\index{monoidal category!symmetric}
  A symmetric monoidal category is a monoidal category $\calC$ together
  with a functorial isomorphism
  $$
    c_{X,Y}\colon X\otimes Y \isoto Y\otimes X
  $$
  such that
  $$
  c_{X,Y} c_{Y,X}=\id
  $$
  and which satisfies the hexagon axiom:
  \begin{center}\includegraphics{figures/fig06.svg}\end{center}
\end{definition}
The hexagon axiom looks complicated; to make it easier to understand, one
can drop all associativity isomorphisms, replacing the hexagon axiom with the
one below
\begin{center}\includegraphics{figures/fig07.svg}\end{center}

\begin{example}\par\noindent
  \begin{enumerate}
    \item Category $\sets$ with respect to disjoint union is a
      symmetric monoidal category; same is true for categories
      $\smoothman$, $\plman$.
    \item
      The category of vector spaces over $\kk$ is a symmetric monoidal category.
      Same is true for categories of representations of a group $G$ or
      a Lie algebra $\g$.
    \item (Super vector spaces). Consider the category $\svect$ whose
      objects are $\Z_2$-graded vector spaces $V=V_0\oplus V_1$. This
      has an obvious monoidal structure.
      Define the commutativity isomorphism $c_{V,W}\colon V\otimes W
      \to W\otimes V$ by
      $$
      c(v\otimes w)=(-1)^{pq} w\otimes v, \quad v\in V_p, w\in W_q.
      $$
      Then this defines on the category $\svect$ a structure of a
      symmetric monoidal category.
  \end{enumerate}
\end{example}
Note that one can also define the symmetric structure on $\svect$ without the
$(-1)^{pq}$ factor; this illustrates the fact that a monoidal category can have
several different symmetric structures.
\begin{remark}
  There are monoidal categories that do not have a natural symmetric structure;
  for example, the category $\vect_G$ of $G$-graded vector spaces, or a
  category
  of bimodules over a ring. There are also ``braided'' categories where we do
  have isomorphisms $c_{X,Y}\colon X\otimes Y \to Y\otimes X$, but these
  isomorphisms fail to satisfy condition $c^2=\id$. The  most famous example of
  a braided  category  is the category of representations of a quantum group. We
  will discuss braided categories later (see \chref{c:321}).
\end{remark}

As before, the real importance of the hexagon axiom is the following coherence
theorem, which is an analog of Mac Lane's coherence theorem, but for
\textbf{symmetric} monoidal categories.

\begin{theorem}[Mac Lane's coherence theorem for symmetric monoidal
 categories]\label{t:maclane-symmetric}\par

  \noindent
  Let $\calC$ be a symmetric monoidal category. Let $F_1, F_2\colon \calC^n\to
  \calC$ be two functors obtained by compositions of
  \begin{itemize}
    \item tensor product functor $\otimes\colon \calC\times \calC\to \calC$
    \item functors $X\mapsto X\otimes \one$, $X\mapsto \one \otimes X$
    \item permutation functor $\calC\times \calC \to \calC\times \calC
          \colon (A, B)\mapsto (B,A)$
  \end{itemize}
  e.g.,
  \begin{align*}
    F_1(A,B,C)&=(B\otimes A)\otimes (\one\otimes C)\\
    F_2(A,B,C)&=C\otimes ((B\otimes  A)\otimes \one).
  \end{align*}
  Then
  \begin{enumerate}
    \item There exists a functorial isomorphism $F_1\simeq F_2$, obtained
      by composition of associativity isomorphism $\al$, left and right
      unit isomorphisms $l,r$, commutativity isomorphism $c$, and their
      inverses.
    \item Any two such isomorphisms $F_1\to F_2$ are equal.
  \end{enumerate}
\end{theorem}
As before, this theorem can be understood as a statement that a certain cell
complex is simply connected; if we only use permutations and tensor product
(i.e., no units), then this complex is called \emph{permutoassociahedron}. (It
took me 3 tries to spell that correctly; in the end, I had to use copy and
paste from another source.)

\clearpage
